\chapter*{Lời cảm ơn}
\label{thanks}

% Tôi xin chân thành cảm ơn ...
% Các bên hỗ trợ API, chart free, 2 thầy hướng dẫn, 

Trước hết, nhóm thực hiện đề tài xin được bày tỏ lòng biết ơn chân thành và sâu sắc đến ThS. Phạm Nguyễn Sơn Tùng và ThS. Hồ Tuấn Thanh, hai giảng viên hướng dẫn đã trực tiếp định hướng, góp ý và đồng hành cùng nhóm trong suốt quá trình thực hiện đồ án tốt nghiệp. Những chỉ dẫn của quý Thầy không chỉ giúp nhóm hoàn thiện hơn về mặt chuyên môn, mà còn giúp nhóm hình thành cách tiếp cận hệ thống hơn trong việc phân tích bài toán, lựa chọn giải pháp kỹ thuật, thiết kế kiến trúc và trình bày kết quả nghiên cứu một cách rõ ràng, khoa học.

Nhóm xin chân thành cảm ơn Ban Giám hiệu Trường Đại học Khoa học Tự nhiên, Đại học Quốc gia Thành phố Hồ Chí Minh, cùng toàn thể quý Thầy, Cô thuộc Khoa Công nghệ Thông tin đã tạo điều kiện học tập, rèn luyện và nghiên cứu thuận lợi trong suốt thời gian qua. Những kiến thức nền tảng về lập trình, cơ sở dữ liệu, công nghệ phần mềm, hệ thống phân tán, an toàn thông tin và phương pháp nghiên cứu được tích lũy trong quá trình học tập là cơ sở quan trọng để nhóm có thể triển khai đề tài ``Xây dựng hệ thống phân tích dữ liệu on-chain Solana và trực quan hóa hành vi ví blockchain''.

Nhóm cũng xin ghi nhận vai trò của các nền tảng, thư viện mã nguồn mở và dịch vụ bên thứ ba đã hỗ trợ quá trình nghiên cứu, phát triển và thử nghiệm hệ thống. Các nguồn dữ liệu và dịch vụ như Helius, Moralis, CoinGecko, Birdeye, Google Gemini, Stripe và hạ tầng RPC của Solana đã cung cấp những cơ sở quan trọng để nhóm có thể tiếp cận dữ liệu on-chain, dữ liệu thị trường, chức năng phân tích hỗ trợ bởi trí tuệ nhân tạo cũng như các luồng thanh toán thử nghiệm trong phạm vi đồ án. Bên cạnh đó, hệ sinh thái mã nguồn mở với các công nghệ như React, Hono, PostgreSQL, Drizzle ORM, ECharts và Carbon Design System cũng đóng vai trò quan trọng trong việc giúp nhóm xây dựng, kiểm thử và hoàn thiện sản phẩm một cách hiệu quả hơn.

Bên cạnh đó, nhóm cũng xin gửi lời cảm ơn đến gia đình và bạn bè đã luôn quan tâm, động viên và tạo điều kiện để các thành viên có thể tập trung hoàn thành đồ án. Sự hỗ trợ về tinh thần, sự cảm thông trong những giai đoạn áp lực và niềm tin từ những người thân xung quanh là nguồn động lực quan trọng giúp nhóm duy trì tinh thần làm việc nghiêm túc trong suốt quá trình thực hiện đề tài.

Cuối cùng, nhóm xin cảm ơn sự nỗ lực, trách nhiệm và tinh thần hợp tác của tất cả các thành viên trong nhóm: Dương Hoàng Hồng Phúc, Lê Võ Minh Phương, Nguyễn Mạnh Phương, Lê Minh Quân, Võ Hoàng Anh Quân và Nguyễn Minh Thiện. Trong quá trình thực hiện đồ án, mỗi thành viên đã cùng nhau trao đổi, phân chia công việc, giải quyết các khó khăn kỹ thuật và hoàn thiện sản phẩm theo từng giai đoạn. Chính sự phối hợp và hỗ trợ lẫn nhau giữa các thành viên là yếu tố quan trọng giúp nhóm hoàn thành đề tài này.

Mặc dù nhóm đã cố gắng thực hiện đề tài một cách nghiêm túc và cẩn trọng, báo cáo khó tránh khỏi những thiếu sót nhất định. Nhóm rất mong nhận được những ý kiến đóng góp quý báu từ quý Thầy, Cô và Hội đồng để có thể tiếp tục hoàn thiện sản phẩm cũng như rút ra thêm nhiều kinh nghiệm cho quá trình học tập, nghiên cứu và phát triển phần mềm trong tương lai.
